% EM structure and teaching revision 20260928
\subsection{有限直导线的积分与端角约定}
令导线沿 $z$ 轴延伸，观察点到轴的垂距为 $a>0$。电流元到观察点的距离为 $\sqrt{a^2+z^2}$，各元产生的磁场方向相同，故
\[
 B=\frac{\mu_0I}{4\pi}\int_{z_1}^{z_2}\frac{a\dd z}{(a^2+z^2)^{3/2}}
 =\frac{\mu_0I}{4\pi a}\left[\frac z{\sqrt{a^2+z^2}}\right]_{z_1}^{z_2}.
\]
若端角从观察点到垂足的方向量起，$z/\sqrt{a^2+z^2}$ 对应带符号的正弦；若从导线方向量起，则应换为余弦。不同教材的端角公式看似不同，往往只是角度定义不同。
令两端趋向正负无穷，结果为 $\mu_0I/(2\pi a)$；令一端在垂足、另一端趋向无穷，结果减半。观察点位于垂足以外时，两端贡献是差而不是两个无符号正弦的和。
对折线或圆弧回路，各段都必须对同一个观察点算矢量再叠加。导线与源点到观察点的方向平行时，该段在毕奥--萨伐尔积分中的贡献为零，但这不意味着该段电流本身不存在。

